\scanpage{24}
\margin{\unsure{Conjectures} Weil, Critères d'Équivalences}

\uline{Théorie de l'équivalence d'Albanese des cycles}

1. \uline{L'anneau des \inserted{\unsure{classes|classe} d'}équivalence d'Albanese}

\uline{Définition} 1.1. Soit $X$ un schéma \deleted{[projectif]} lisse sur le corps $k$, $\unsure{\mathfrak z|z}$ un cycle sur $X$. On dit que $\unsure{\mathfrak z|z}$ est alb-équi. à $0$, si on peut trouver un schéma [projectif] lisse \deleted{\unsure{géométriquement}} \inserted{connexe} $\unsure{T|Y}$ sur $k$, \deleted{et} un $0$-cycle \deleted{de degré $0$ \unsure{$\nu$}} $\nu$ sur $\unsure{T|Y}$, et un cycle $Z$ sur $X\times\unsure{T|Y}$, tel que
\[
\text{(i) }\inserted{\nu\text{ est de degré }0\text{ et }}S_Y(\nu)=0\qquad\text{(ii) }Z(\nu)\text{ est défini et}
\]
égal à $\unsure{\mathfrak z|z}$.

\margin{?}

\uline{Proposition} 1.2. (i) Les cycles alb-équi. à $0$ sur $X$ forment un sous-groupe \inserted{gradué $N^*(X)$,} de $\operatorname{Cyc}^*(X)$, contenant le groupe des cycles lin. équiv. à $0$, de sorte que la connaissance de $N^*(X)$ équivaut à celle de son image $N'^*(X)$ dans \deleted{\unsure{$\operatorname{Cyc}$}} $A^*(X)$ (l'anneau de Chow de $X$, défini pr les classes à équiv. linéaire près), et $\operatorname{Cyc}(X)/N(X)\simeq\deleted{\unsure{\operatorname{Cyc}}}A^*(X)/N'^*(X)$.

\mspage{23}
